\lecture{16}{Mon 17 Oct 2022 10:38}{}

\begin{lemma}
	(2.9.3) If $G$ is an internal direct product of its normal subgroups $N_1, \ldots, N_n$, then $N_i \cap N_j = \{e\} $ for $i \neq j$.
\end{lemma}
\begin{proof}
	If $a \in N_i \cap N_j$ and $a \neq  e$ then
	\[
		a = e_1e_2\ldots a_i e_n\ldots e 
		= e_1 e_2 \ldots a_j \ldots e_n
	.\] 
	which contradicts the uniqueness of representation of $a$ in the direct product.
\end{proof}

\begin{lemma}
	(2.9.2)
	Let $G$ be a group with $M \triangleleft G$ and $N \triangleleft G$ and $M \cap  N  =\{e\} $. Then for all $m \in  M$ and $n \in N$ one has $mn = nm$.

\end{lemma}
\begin{proof}
	Observe that if $a = m n m^{-1} n^{-1} $, then $a = m(n m^{-1} n^{-1} ) \in M$. Also, $a = (m n m^{-1} ) n \in N$. That means $a \in M \cap N$, so $a = e$. Thus $mn = nm$.
\end{proof}
\begin{theorem}
	(2.9.4)
	Let $G$ be a group with $N_1,\ldots,N_n$. Then the map
	\begin{align*}
		\psi: N_1\times N_2\times \ldots\times N_n &\longrightarrow G \\
		(g_1,g_2,\ldots,g_n) &\longmapsto \psi((g_1,g_2,\ldots,g_n)) = g_1g_2\ldots g_n
	.\end{align*}
	is an isomorphism if and only if $G$ is the internal direct product of $N_1, \ldots, N_n$.
\end{theorem}
\begin{proof}
	(Proof of 2.9.4)
	
	Suppose $\psi$ defined in hypothesis is an isomorphism. Then each $g \in G$ has a representation $g = g_1g_2\ldots g_n$ for some $g_i \in N_i$ by surjectivity. Also by injectivity, whenever $g = g_1\ldots g_n$ with $(g_1, g_2, \ldots, g_n) \in N_1\times \ldots \times N_n$, each $g_i$ is unique, so representation is unique. Then $G$ is the internal direct product of $N_1, \ldots, N_n$.

	Suppose $G$ is the internal direct product of $N_1, \ldots, N_n$. Then each $g \in G$ has a unique representation $g = g_1g_2 \ldots g_n$. So $\psi$ is surjective. If $\psi(g_1, \ldots, g_n) = \psi(h_1, \ldots, h_n)$, then uniqueness of representation shows $g_1 \ldots g_n = h_1 \ldots h_n$ implies  $g_i = h_i$ for each $i$. So the $n$-tuples are the same. So $\psi$ is injective.

	Remains to show $\psi$ is a homomorphism. But whenever $(g_1, \ldots, g_n)$ and $(h_1, \ldots, h_n)$ lie in $N_1 \times  \ldots \times  N_n$, then
	\begin{align*}
		&\psi\left( (g_1, \ldots, g_n)(h_1, \ldots, h_n) \right)\\
		&= \psi(g_1h_1,\ldots ,g_n h_n) \\
		&= g_1h_1 g_2h_2 \ldots g_n h_n \\
		&= \left( g_1g_2\ldots g_n \right)  \left( h_1 h_2 \ldots h_n \right)  & \text{Lemma 2.9.2}\\
		&= \psi(g_1,\ldots,g_n)\psi(h_1,\ldots,h_n) 
	\end{align*}
	Thus $\psi$ is a homomorphism, it's injective and surjective, hence it is an isomorphism.
\end{proof}

\subsection{Finite Abelian Groups}
\begin{theorem}[Classification for Finite Abelian Groups]
	(2.10 A)
	Suppose that $G$ is a finite abelian group. Then $G$ is isomorphic to a direct product of cyclic groups of the shape \[
		\mathbb{Z}_{p_1^{\alpha_1}}\times \mathbb{Z}_{p_2^{\alpha_2}}\times\ldots \times  \mathbb{Z}_{p_r^{\alpha_r}}
	.\] 
	where $p_i$ are primes (not necessarily distinct), and $\alpha_i \in \mathbb{N}$, and $n = |G| = p_1^{\alpha_1} \ldots p_r^{\alpha_r}$.
\end{theorem}

\begin{corollary}
	(2.10 B)
	If $G$ is a finite abelian group, then there are integers $m_1, \ldots, m_k$ with $2 \le m_1 \le  m_2 \le  \ldots \le  m_k$, and $m_i | m_{i+1}$ for $1\le i<k$, such that \[
		G \cong \mathbb{Z}_{m_1} \times \mathbb{Z}_{m_2} \times \ldots \times \mathbb{Z}_{m_k}
	.\] 
\end{corollary}

\begin{example}
	Classify all abelian groups of order $1400$.

	Observe that $1400 = 2^3 5^2 7$. Then by classification theorem, if $|G| = 1400$, then $G$ is isomorphic to one of:
	\begin{itemize}
		\item $\mathbb{Z}_2 \times \mathbb{Z}_2 \times \mathbb{Z}_2 \times \mathbb{Z}_5 \times \mathbb{Z}_5 \times \mathbb{Z}_7 \cong \mathbb{Z}_{70} \times \mathbb{Z}_{10} \times \mathbb{Z}_2$
		\item $\mathbb{Z}_2 \times \mathbb{Z}_4 \times \mathbb{Z}_5 \times \mathbb{Z}_5 \times \mathbb{Z}_7 \cong \mathbb{Z}_{140} \times \mathbb{Z}_{10}$ 
		\item $\mathbb{Z}_8 \times \mathbb{Z}_5 \times \mathbb{Z}_5 \times \mathbb{Z}_7 \cong \mathbb{Z}_{280}\times \mathbb{Z}_5$
		\item $\mathbb{Z}_2 \times \mathbb{Z}_2 \times \mathbb{Z}_2 \times \mathbb{Z}_{25} \times \mathbb{Z}_7\cong \mathbb{Z}_{350} \times \mathbb{Z}_2 \times \mathbb{Z}_2$ 
		\item $\mathbb{Z}_2 \times \mathbb{Z}_4 \times \mathbb{Z}_{25} \times \mathbb{Z}_7 \cong \mathbb{Z}_{700} \times \mathbb{Z}_2$
		\item $\mathbb{Z}_8 \times \mathbb{Z}_{25} \times \mathbb{Z}_7 \cong \mathbb{Z}_{1400}$
	\end{itemize}
	There are $6 $ isomorphism classes for all finite abelian group of order $1400$.
\end{example}
